\documentclass[../main.tex]{subfiles}
\begin{document}
%==========================================TIÊU ĐỀ TẬP========================================
\begin{table}[h]
  \begin{adjustwidth}{-1cm}{}
    \begin{tabular}{>{\centering\arraybackslash}p{6cm}|>{\raggedright\arraybackslash}p{12.5cm}}
      \multirow{5}{6cm}
      % Cột bên trái: ÉPISODE và số tập
      {\\[-40pt]
        \flaregothic\fontsize{20pt}{20pt}\selectfont \centering
        {\contour{errorthree}{\textcolor{black}{ÉPISODE}}}\color{black}\\
        \burbank\fontsize{72pt}{72pt}\selectfont
        \includegraphics[height=3cm]{../../../../Image/Book?/number/num10.png}
        \\[18pt]
        % \flaregothic\fontsize{14pt}{14pt}\selectfont
        % \color{epattr}BIENVENUE, \uline{\color{Banpire} Mel} !
        % \flaregothic\fontsize{18pt}{18pt}\selectfont
        % \color{epattr}ÉPISODE PERDU
        \color{black}
      }
       &
      % Dòng thứ nhất: tên tập tiếng Nhật
      {\fotskip\fontsize{18pt}{18pt}\selectfont \ruby{日}{まい}\ruby{常}{にち}\ruby{的}{の}ERROR
      }                                         \\[6pt]
      % Dòng thứ hai: tên tập tiếng Anh
       & {\cafeteria\fontsize{18pt}{18pt}\selectfont ERROR in Normalcy}                                                         \\[4pt]
      % Dòng thứ ba: tên tập tiếng Việt
       & {\vnmsans\fontsize{18pt}{18pt}\selectfont Sự bất thường trong sự bình thường}                                          \\[8pt]
      % Dòng thứ tư: Nhân vật xuất hiện
       & {\caviar{\fontsize{14pt}{14pt}\selectfont Track used: The Hollowness (from Dancing Line: Community Edition/Max Line)}}
      \\[6pt]
      % Dòng cuối cùng: Thời gian diễn ra của tập (thường sẽ là ngày phát hành)
       & {\continuum\fontsize{16pt}{16pt}\selectfont Ngày phát hành: 22/04/2022}                                                \\[8pt]
    \end{tabular}
  \end{adjustwidth}
\end{table}
%===========================================NỘI DUNG==========================================
\color{black}

Đêm phủ kín nhà ga Aogami.

Khung cảnh vốn đông đúc và ồn ào ban sáng giờ biến thành một khoảng lặng bất thường. Không còn tiếng giày dép lẫn tiếng loa phát thanh quen thuộc, chỉ còn khoảng không mênh mông tối đen bám chặt lấy từng vách tường, từng tấm bảng chỉ dẫn. Những bóng đèn huỳnh quang treo dọc trần vẫn sáng, tỏa ánh sáng trắng nhợt nhạt, nhưng không xua nổi cái tĩnh lặng rợn ngợp đang vây quanh.

Sân ga trống rỗng, gợi cảm giác như nơi này đã bị bỏ quên hàng thập kỷ. Thế nhưng, giữa khung cảnh im lìm ấy, một bóng hình khổng lồ vẫn hiện diện: đoàn tàu lạ lẫm, thân sắt đen kịt phủ lớp bụi mờ, đang đứng chờ bất động trên đường ray. Trên đầu máy, một tấm biển kim loại mờ xỉn ánh sáng phản chiếu: 211493. Những con số như khắc sâu vào bóng đêm, vừa vô nghĩa, vừa gợi cảm giác như một mật mã chết chóc.

\img{e2-2a1}

Cửa toa mở sẵn, hé ra khe hở tối om, như thể đang mời gọi.

Nekomiya Yuka đứng trước cửa toa, đôi mắt hiếu kỳ nhìn vào bóng tối bên trong. Bỗng nhiên, từ trong khoang tối đen như mực ấy, một bàn tay mảnh khảnh, trắng bệch như sương giá, từ từ vươn ra, ngón tay thon dài chỉ về phía cô, như một lời mời gọi câm lặng. Không ai có thể nhìn thấy chủ nhân của bàn tay ấy, nhưng Yuka dường như hiểu được điều gì đó mà không ai khác có thể cảm nhận.

Cô bước vào, không một chút do dự, không một lời gọi, chẳng hề ngoái nhìn lại phía sau. Cô đi như kẻ bị thôi miên, bước chân đều đặn vang lên nhè nhẹ trên nền ga, hòa vào tiếng kim loại lạnh buốt vọng lại từ khung toa tàu. Bàn tay cô lướt qua mép cửa, rồi cả thân người biến mất vào trong, như thể đó là nơi cô vốn dĩ phải đến. Trong khoảnh khắc cuối cùng trước khi hoàn toàn biến mất, dường như có một nụ cười nhẹ thoáng qua trên môi cô, như thể cô đang đón nhận một lời mời từ ai đó mà chỉ mình cô nhìn thấy.

\img{e2-2a2}

*XOẠCH*

Ngay khi cánh cửa khép lại, đoàn tàu run lên, kim loại rít vang như tiếng thở dài. Khói trắng phụt ra từ ống khói, tản vào màn đêm như một lớp sương mù nặng trĩu. Tiếng còi tàu chát chúa xé toạc không gian, kéo dài thành âm hưởng như từ thời đại nào xa lắc vọng về.

*\textbf{TU TU!}*

Bánh sắt bắt đầu lăn.

Nhịp kim loại đều đặn, rền rĩ, dần tăng tốc. Ánh đèn toa tàu sáng mờ xuyên qua những ô cửa phủ bụi, như hàng chấm sáng lạc lõng chạy dọc đường ray. Trong khoảnh khắc ngắn ngủi ấy, bóng dáng đoàn tàu hiện lên đầy bí ẩn, rồi chỉ ít giây sau, cả thân tàu chìm hẳn vào màn đêm, tan biến khỏi sân ga, để lại khoảng không trống rỗng đến lạnh người.

Lặng lẽ, nơi cửa sân ga, một bóng người vẫn đứng đó từ trước: Misora Shino.

Không một lời nói. Không một cử chỉ. Đôi mắt xanh của cô ánh lên thứ gì đó khó tả, không phải buồn, không phải vui, cũng chẳng phải tiếc nuối, mà giống như đang chứng giám cho một điều hiển nhiên mà sức cô không thể ngăn cản. Ánh sáng trắng nhợt từ đèn trên cao phủ xuống gương mặt ấy, làm cho vẻ tĩnh lặng càng trở nên xa xăm.

Rồi, cũng như chính khung cảnh vô thực này, Shino dần nhạt nhòa, biến mất, để lại sân ga trống rỗng.

Chỉ còn màn đêm thanh tĩnh.

Chỉ còn tiếng vang của một chuyến tàu đã đi xa, xa đến mức không ai biết nó rốt cuộc sẽ về đâu.

\pov{Haragen Saikyou}

\dots

\dots

``!?''

Tôi giật bắn mình, tỉnh dậy ngay trên giường.

Một luồng hơi lạnh chạy dọc sống lưng khi đôi mắt mở to ra khỏi bóng tối nặng nề của giấc mơ vừa rồi. Cổ họng khô khốc, nhịp tim đập thình thịch như vừa chạy ma-ra-tông. Tôi ngồi bật dậy, bàn tay vô thức nắm chặt lấy chăn, trong đầu vẫn còn in hằn những âm thanh rền vang của đoàn tàu kia, thứ âm hưởng như đang lẩn khuất đâu đó trong thính giác tôi, dù xung quanh tôi lúc này chỉ có sự thanh bình của buổi sáng.

Khẽ quay sang bên, tôi sững lại. Nii-chan và Reiko đã tỉnh từ bao giờ. Anh tôi ngồi sát cửa sổ, tay chống cằm, ánh nắng buổi sáng chiếu nghiêng trên gương mặt cậu, tay cầm cốc cà phê, còn Reiko thì đang xếp gọn bộ quần áo của mình lại.

Tôi buột miệng hỏi, giọng khàn hẳn: ``Hai người\dots\ dậy từ khi nào vậy?''

Nii-chan không quay lại, chỉ nhàn nhạt đáp: ``Năm giờ sáng.''

Tôi há hốc miệng, suýt nghẹn luôn hơi thở. ``S-sớm thế!?''

Anh chỉ nhún vai, môi khẽ nhếch thành nụ cười nửa miệng: ``Xin lỗi. Do thói quen thôi.

Reiko bật cười khúc khích, rồi dịu giọng: ``Em thì vừa mới dậy thôi. Nhưng mà, Onee-sama\dots\ sao vừa tỉnh dậy mà giật mình ghê vậy?''

Tôi chần chừ một thoáng, rồi siết chăn lại. Cái cảm giác ấy còn quá rõ rệt để mà lảng tránh.

``Chị\dots\ Chị vừa có một cơn ác mộng.''

Cả hai cùng hướng ánh nhìn về phía tôi. Nii-chan im lặng, Reiko nghiêng đầu lo lắng. Tôi thở dài, quyết định nói ra hết.

``Trong mơ, chị thấy mình ở nhà ga Aogami. Nhưng\dots\ không phải buổi sáng như hôm qua. Là ban đêm. Tối đen, chỉ có ánh đèn nhợt nhạt. Trên sân ga có một con tàu\dots\ mang biển số 211493.'' Tôi dừng lại một nhịp, tim vẫn còn nện dồn dập, rồi tiếp tục: ``Thế rồi, Yuka-san\dots\ chị thấy cô ấy bước lên con tàu đó. Có cảm giác như ai đó mời gọi. Cánh cửa đóng lại, rồi nó rời đi. Biến mất luôn trong bóng đêm.''

Không gian lặng đi một lúc. Em tôi rùng mình, còn Nii-chan thì hơi cau mày. Tôi nhận ra vẻ mặt anh thoáng đanh lại khi nghe đến hai chữ ``nhà ga.''

``Nhà ga\dots\ Ý em là cái nhà ga Aogami đó?'' anh lặp lại, giọng trầm thấp.

``Đúng vậy\dots\ ``tôi gật đầu, cảm giác gai lạnh nổi lên từng đợt. ``Chính nó. Điều mà Kaoru-san đã nói hôm qua\dots\ Ngôi đền Aogami, rồi nhà ga Aogami\dots\ tất cả xoay quanh cái tên đó. Thật sự\dots\ chẳng lẽ chỉ là trùng hợp thôi sao?

Reiko khẽ ôm lấy cánh tay tôi, giọng nhỏ xíu: ``Em cũng thấy bất an. Không biết có liên quan đến mấy lời đồn người mất tích trong thị trấn không nữa\dots''

Tôi rùng mình. Phải rồi. Những lời đồn ấy. Những cái tên biến mất quanh ga. Cái cảm giác quen thuộc quái gở.

``Nii-chan, anh cũng thấy như vậy\dots\ đúng không? ``tôi cất tiếng, nhìn thẳng vào anh trai.

Anh thở dài, gật khẽ.

``Những gì chúng ta thấy, nghe và cảm nhận từ lúc ta rời nhà ga tới giờ\dots\ không thể nhầm lẫn vào đâu được. Quả thực, nhà ga ấy có vấn đề.''

Tôi cắn môi. Câu nói ấy làm tim tôi thắt lại. Thật sự, chúng tôi đâu thể lảng tránh được mãi. Nhưng dù vậy\dots\ hôm nay chúng tôi vẫn phải tiếp tục, phải để mọi thứ diễn ra như ngày bình thường.

Tôi hít sâu lấy bình tĩnh, rồi hỏi: ``Bây giờ là mấy giờ rồi?''

Nii-chan nhìn chiếc điện thoại, trả lời ngắn gọn: ``Tám giờ rưỡi.''

``Cái gì!? ``tôi bật dậy, cuốn chiếc chăn ném sang một bên, cảm giác bối rối thay thế cả nỗi sợ vừa rồi. ``Mau chuẩn bị đi thôi! Chúng ta còn phải gặp Akane-san ở khu phố truyền thống lúc chín giờ!''

Nii-chan khẽ cười, nhưng kèm theo một câu làu bàu: ``Reiko thì chuẩn bị từ thuở nào rồi, còn anh thì nhanh lắm. Còn mỗi em thôi đấy.''

Tôi giơ tay phản bác ngay, mặt đỏ bừng: ``E-em chuẩn bị cũng nhanh mà!''

Nói xong, tôi vội vã chạy vào nhà tắm, tiếng dép loẹt xoẹt vang khắp phòng, để lại tiếng cười nhẹ pha chút bất lực của em tôi và tiếng thở dài quen thuộc của Nii-chan phía sau.

\begin{center}
  \uline{8:50 sáng.}
\end{center}

Chúng tôi đã có mặt ở con phố truyền thống mà Akane-san đã hẹn. Bầu trời hôm nay sáng trong, gió xuân nhẹ thổi qua, mang theo mùi hương bánh nướng từ một tiệm nào đó gần đây.

Thế nhưng, điều làm tôi thấy kì lạ chẳng phải là khung cảnh xung quanh\dots\ mà lại nằm ở chính chúng tôi. Cụ thể hơn, là bộ trang phục tôi đang mặc.

Tôi nhìn lại bản thân mình một lần nữa rồi nhăn mặt: ``Này, sao cả ba chúng ta lại mặc đồng phục trường vậy chứ?''

Tôi còn nhớ rất rõ, đó chính là bộ chúng tôi đã mặc hôm qua. Tôi chắc mẩm rằng ít nhất sáng nay chúng tôi sẽ mặc gì đó thoải mái hơn. Nhưng rõ ràng, cả tôi, Nii-chan và Reiko đều đồng bộ như ba học sinh chuẩn bị vào tiết học - mà cũng chẳng phải, vì hôm nay là ngày nghỉ!

Anh trai tôi chỉ nhún vai, vẻ mặt bình thản đến mức phát bực: ``Lục tung cả nhà thì cũng chỉ có mỗi bộ này với bộ đồ ngủ thôi em.''

``\textbf{Hả!?}'' tôi tròn mắt. ``B-Bố mẹ vẫn chưa chuyển quần áo sang cho chúng ta sao!?''

``Chưa. Anh mở hết các hộp rồi, toàn thấy xoong chảo với mấy đồ lặt vặt thôi.''

``Cái gì chứ\dots'' tôi tặc lưỡi, ngán ngẩm với ông bà già.

Nụ cười nửa miệng của anh ấy xuất hiện ngay sau đó: ``Hay là em muốn mặc luôn đồ ngủ ra ngoài?''

``Đ-đừng có trêu em nữa!'' tôi đỏ mặt, quay đi, dù trong lòng cũng tự nhủ\dots\ đúng là chẳng còn lựa chọn nào khác thật.

Reiko thì lại nhẹ nhàng, kéo tay áo len đồng phục lên cao một chút, miệng cười hồn nhiên: ``Em thấy mặc thế này cũng ổn mà. Lại còn ấm nữa.''

\dots Tôi không hiểu nổi tại sao Nii-chan và Reiko không cảm thấy ngại khi diện bộ đồ có vẻ lạc quẻ đó. Hay là, chắc chỉ mình tôi là cảm thấy không thoải mái. Thôi thì, ít nhất thì còn hơn phải mặc bộ đồ ngủ nhếch nhác ấy ra khỏi nhà.

Chúng tôi đứng chờ ở đầu phố. Người qua kẻ lại khá đông, nhiều cửa tiệm đang mở hàng. Nii-chan khoanh tay, thi thoảng liếc nhìn đồng hồ trên điện thoại. Reiko thì nghịch mấy dải giấy omikuji\footnote{Phong tục bói toán truyền thống ở Nhật Bản, nơi mọi người rút một quẻ giấy tại các ngôi đền hoặc chùa để nhận dự đoán về vận mệnh, vận may, hoặc lời khuyên cho các khía cạnh như tình yêu, sức khỏe, công việc và học tập.} buộc ở một góc gần đó, trông vô tư hơn hẳn. Còn tôi\dots\ thì cứ đứng nhìn dòng người, và rồi ánh mắt vô thức dừng lại ở phía xa con đường kéo dài về ga Aogami.

Hình ảnh trong giấc mơ ban nãy lại lởn vởn quay về. Con tàu mang số hiệu 211493. Yuka-san bước vào, có cảm giác như ai đó mời gọi, biến mất không dấu vết. Rồi ánh mắt trống rỗng thể hiện sự bất lực của Shino-san đứng từ xa dõi theo.

Mọi thứ nhìn bề ngoài có vẻ rời rạc, nhưng trông chúng có vẻ có một mối liên hệ mơ hồ nào đó. Những lời đồn về người mất tích. Lớp học vắng đi từng bạn một. Trận cúm mùa mà Kaoru-san đã nhắc. Rồi cả lời Shino-san hôm qua\dots

Âm vang ấy vẫn như vang vọng bên tai. Tôi siết chặt bàn tay mình, cố ngăn không cho ý nghĩ đáng sợ đó nuốt chửng mình thêm lần nữa.

Anh tôi liếc sang, như đoán được suy nghĩ, nhưng không nói gì, chỉ lặng lẽ đứng cạnh như một sự trấn an thầm lặng.

Có lẽ\dots\ một lúc nào đó, chúng tôi có thể khám phá ra được bí ẩn đằng sau ga Aogami kia.

\ct

Chúng tôi chờ thêm mười phút. Và rồi, ở đầu phố, bóng dáng ba cô gái xuất hiện, dáng vẻ họ vừa đi vừa trò chuyện rôm rả, tiếng cười xen lẫn với không khí nhộn nhịp buổi sáng.

Tôi nheo mắt nhìn về phía đó, dò xét một hồi.

\dots

\dots

Chính là họ. Saya-san, Kaoru-san và Miku-san.

Vẫn là cô gái cầm hộp sữa ấy, vẫn là cô gái tóc vàng kia, và vẫn là cô gái có bờm sói cầm thanh sô-cô-la đó.

\img{e2-2c}

``Chào buổi sáng ba người\~{}'' Saya-san lên tiếng đầu tiên, mắt vẫn nhắm hờ.

``Các cậu đã chờ lâu chưa?'' Miku-san hỏi, tay cầm thanh sô-cô-la ăn dở, giọng nhẹ như gió thoảng.

Kyuen khẽ nghiêng đầu, đáp bằng cái giọng lịch sự thường ngày: ``Bọn tôi cũng vừa mới tới thôi.''

Tôi liếc nhìn quanh rồi khẽ cau mày. Có gì đó\dots\ thiếu thiếu từ nhóm bạn. Akane-san và Yuka-san đang không có ở đây. Chính họ đã chủ động hẹn gặp chúng tôi hôm nay mà nhỉ?

Tôi mở miệng: ``Ơ\dots\ còn Akane-san và Yuka-san thì--''

Lời chưa dứt, một cơn chấn động kì lạ lướt qua mắt tôi. Cảnh vật như bị phủ một lớp nhiễu mờ, các hình ảnh trước mặt chập chờn như đoạn băng cũ đang phát lại.

``C-Cái quái gì đang xảy ra vậy\dots?''

Và rồi, Miku quay sang hỏi: ``Có địa điểm nào cụ thể mà các cậu muốn tham quan không?''

Tôi giật mình.

Đó chính là\dots\ câu hỏi mà cô ấy đã hỏi chúng tôi hôm qua. Giống hệt. Không sai một chữ.

``Ơ\dots\ chẳng phải cậu đã hỏi bọn tôi như thế sao?'' tôi buột miệng, giọng run nhẹ.

Nii-chan cũng chau mày, tiếp lời: ``Có chứ, nhưng\dots\ còn Akane-san và Yuka-san thì sao?''

Saya-san bỗng chen vào, lời nói vang lên rành rọt, từng âm như in đậm trong đầu tôi: ``Nếu vậy thì các cậu thử đến ngôi đền Aogami xem nào?''

Tôi chết lặng. Một tia sét đánh ngang tai.

Đó cũng chính là câu mà cô ấy đã nói hôm qua. Đúng từng từ một, không lệch một từ.

Chưa kịp để tôi phản ứng, Kaoru nâng hộp sữa trong tay, như một động tác quen thuộc: ``Đó cũng là tên của ngôi đền vùng này.''

Từng sợi lông tơ sau gáy tôi dựng ngược.

Không. Đây không phải trò đùa. Họ đang lặp lại những gì họ nói. Y như một con rô-bốt.

Tôi đã từng nhìn thấy cảnh tượng này ở đâu rồi.

Hình ảnh thoáng vụt qua trong đầu tôi: Hanabi và Inari, sau trận động đất, cũng từng lặp lại lời nói của chính mình, một cách vô hồn và rập khuôn.

Và để khẳng định nỗi kinh hoàng đang dâng lên, Kaoru-san bất chợt nói thêm, giọng dõng dạc đầy tự tin: ``Tất nhiên rồi! Đã thế nha, còn chung lớp nữa chứ.''

Nhưng\dots\ không ai trong chúng tôi hỏi câu đó cả.

Tôi đứng chôn chân tại chỗ, bàn tay siết chặt, tim đập dồn dập, đầu óc tôi bắt đầu loạn dần, không thể nào phân biệt nổi đâu là thực, và đâu là ảo giác.

Nii-chan cũng cảm nhận thấy sự bất thường, đôi mắt cậu ánh lên vẻ cảnh giác. Anh không dám ho he một tiếng, mặc dù trong thâm tâm anh cũng muốn hỏi họ vài điều.

Chỉ có Reiko ngơ ngác nhìn hết người này sang người khác, môi mấp máy: ``Ể\dots? Sao mọi người\dots\ lặp lại y nguyên như vậy? Mình tưởng\dots\ chúng ta đã thống nhất với nhau từ trước rồi chứ?''

Không khí quanh chúng tôi bỗng trở nên lạnh lẽo đến rợn người, như thể có một bức màn vô hình đang bao phủ lấy nhóm. Cái lạnh ấy thật phi lý và trái ngược - bởi lẽ trên cao kia, vầng dương vẫn đang tỏa nắng rực rỡ, những tia nắng vàng ấm áp vẫn đang chiếu rọi xuống con phố, nhưng dường như chúng không thể xuyên thấu được bầu không khí kỳ quái đang vây lấy chúng tôi.

Reiko lên tiếng trước, giọng nhỏ nhẹ nhưng có chút sốt ruột: ``Thế\dots\ Akane-san và Yuka-san đâu rồi?''

Tôi chờ đợi một màn lặp lại máy móc như lúc nãy, nhưng lần này\dots\ họ trả lời khác.

Miku-san mỉm cười, gương mặt tươi tắn quen thuộc: ``Akane-san sẽ tới đây sớm thôi.''

Saya-san tiếp lời, tông giọng đầy hồn nhiên: ``Ừm! Mình cũng nghĩ vậy đấy\~{}''

Ngay khoảnh khắc ấy, tôi thở phào, đôi vai nhẹ bẫng như vừa được dỡ  xuống một tảng đá. Không chỉ tôi, mà Nii-chan cũng khẽ gật đầu, còn Reiko thì buông ra một tiếng ``À\dots'' đầy nhẹ nhõm. Ít ra, họ không còn lặp lại như rô-bốt nữa. Tôi tự nhủ: ``Có lẽ mình đã tưởng tượng quá mức\dots\ Có lẽ mọi chuyện vẫn ổn.''

Nhưng rồi, giọng của Kaoru-san phá tan cái vỏ an toàn mỏng manh đó: ``Mà này, các cậu\dots\ \emph{Yuka là ai vậy?}''

Câu hỏi hồn nhiên ấy rơi xuống, lạnh lẽo như một tảng băng. Cảm giác như thể một cơn gió lạnh thổi qua mạnh tới độ làm cho tôi đông cứng.

Đôi mắt của ba người họ lập tức đổi khác. Ánh sáng trong đó vụt tắt, thay bằng một lớp đục mờ, vô hồn như mắt cá chết. Tôi từng thấy ánh mắt đó. Vẫn là cô bạn tóc đuôi ngựa buộc lệch, cô bạn bờm cáo trắng ấy\dots\ sau trận động đất, họ cũng đã nhìn tôi như thế.

\img{e2-2d}

Nii-chan cũng không giấu nổi sự bất ngờ, liền lắp bắp, giọng khô khốc: ``H-Hả? Mọi người đang nói gì vậy?''

Reiko thì bật lên phản ứng bản năng, đầy phản kháng: ``Ba cậu hỏi gì kỳ quặc vậy? Cô gái có bờm mèo, tóc tím ngắn, thân thiết với Akane-san ấy? Chính là Yuka-san đó!''

Ba cô gái nghe tới đó chỉ nhìn nhau. Họ cùng đồng loạt nghiêng đầu, động tác chậm chạp, như một vở kịch diễn sai nhịp.

Cô gái tóc vàng mở lời, giọng chậm rãi và kéo dài, nhưng trong đó có thứ gì đó không phải là con người: ``Miku-san\dots\ có phải người đó, là bạn cậu không?''

Cô bạn tai sói chớp mắt, rồi đáp thẳng một chữ lạnh tanh: ``Không\dots''

Cô quay sang hỏi ngược lại, cũng bằng giọng trống rỗng: ``Bạn cậu à?''

Saya lắc đầu.

Không khí trở nên nặng nề đến nghẹt thở. Tôi thấy tay mình run lên, dù đã cố nắm chặt vạt áo. Trong đầu, hình ảnh cơn ác mộng từ sáng nay hiện về: Yuka, đứng dưới ánh đèn ga tàu, có cảm giác như ai đó mời gọi, đã lên chuyến tàu đó, rồi biến mất vào bóng tối.

Tôi nuốt khan. Anh trai tôi quay sang nhìn tôi, ánh mắt hoảng loạn, còn Reiko thì bấu chặt cánh tay tôi. Chúng tôi cùng lúc rùng mình, bởi liên tưởng đó trùng khớp với những lời đồn dai dẳng quanh nhà ga: những người bị cuốn vào vòng xoáy, rồi biến mất\dots\ không bao giờ trở lại.

Trong tâm trí ba chúng tôi, câu hỏi ấy vang dội rõ ràng: ``Không lẽ\dots\ Yuka-san đã--!?''

Kaoru-san bất chợt cắt ngang mạch suy nghĩ hỗn loạn của chúng tôi. Cánh tay cô ấy đưa về phía tôi, nụ cười rạng rỡ bừng sáng trên gương mặt, như thể chưa từng có chuyện gì xảy ra.

``Mà, thôi kệ đi,'' cô cất giọng vui tươi. ``Để bọn tôi dẫn cậu tham quan cho.''

Ngay sau đó, Saya cũng nghiêng đầu, đôi mắt long lanh như cười: ``Phải đấy! Đi chung sẽ vui hơn nhiều.''

Miku chắp tay sau lưng, giọng ngọt lịm: ``Bọn tôi còn biết kha khá chỗ hay ho quanh đây nữa. Nếu các cậu không phiền, tôi có thể chỉ luôn.''

\img{e2-2e}

Tôi và Nii-chan đồng loạt rùng mình. Nụ cười kia quá sáng, quá trọn vẹn, đến mức không còn là nụ cười của con người. Nó giống như lớp mặt nạ được đắp lên khuôn mặt rỗng tuếch bên dưới. Tôi cảm nhận được, một thế lực nào đó đang điều khiển họ, muốn kéo chúng tôi rời xa khỏi thứ sự thật kinh hoàng đang bị che giấu.

Chúng tôi đồng loạt lùi lại một bước. Tôi run rẩy, mở miệng lắp bắp: ``X-Xin lỗi các cậu\dots\ Tôi không thể đi được rồi.''

Reiko tiếp lời, tỏ vẻ áy náy: ``Bọn mình cần phải tìm Yuka-san trước\dots\ Và cả Akane-san nữa. Cho nên\dots\ để hôm khác nhé?''

Anh trai tôi thì nghiêm mặt, không rườm rà, nhanh chóng hối thúc: ``Đi thôi, Saikyou, Reiko. Nhanh lên.''

Bản năng mách bảo chúng tôi: nơi cần đến chính là nhà ga Aogami. Không cần bàn bạc, cả ba định quay lưng chạy về hướng đó.

Nhưng bất ngờ thay, Kaoru-san vươn tay chộp lấy cổ tay tôi, siết chặt. Nụ cười tắt ngấm, thay bằng giọng nói trầm thấp, đáng sợ với nét mặt nghiêm lại: ``Đừng. Đừng đi về phía đó.''

\img{e2-2f}

Cô bạn tóc vàng tiến lên một bước, đôi mắt vô hồn nhìn xoáy vào tôi: ``Không ai nên tới ga vào lúc này đâu\~{}''

Còn cô bạn tai sói khẽ lắc đầu, lời thì thầm lạnh lẽo vang lên bên tai tôi: ``\dots Nếu đi, các cậu sẽ hối hận đấy.''

Tôi rùng mình, cảm thấy tức ở lồng ngực. Cơn sợ hãi biến thành sức mạnh bất ngờ, tôi giật mạnh tay ra, hét lên: ``\textbf{Bỏ tay tôi ra!}''

Rồi không để lại một giây do dự, tôi cùng Nii-chan và Reiko lao đi về phía ngược lại, bỏ lại sau lưng ba nụ cười vô hồn kia, đứng bất động như những bức tượng gãy nhịp trong màn kịch.

\ct

Chúng tôi dừng lại ở một công viên, nơi những tán cây anh đào và phong lá đỏ đổ bóng che phủ một khoảng đất rộng. Nhìn bề ngoài, nơi này vẫn mang dáng vẻ của một công viên bình dị: những chiếc ghế đá xám nằm rải rác dưới tán cây, lối đi lát gạch uốn lượn như dải lụa, và đám cỏ xanh mơn mởn trải dài tới tận chân trời. Thế nhưng, chúng tôi không còn dám tin tưởng vào vẻ yên bình đó nữa. Bởi lẽ, ranh giới giữa thực và ảo đã trở nên mong manh đến độ, ngay cả không gian tưởng chừng an toàn này cũng có thể chỉ là một ảo ảnh khác của thị trấn kỳ lạ này mà thôi.

Cả ba cùng ngồi phịch xuống ghế đá, hơi thở gấp gáp như vừa chạy một quãng đường dài. Tôi đưa tay lên ngực, cảm nhận nhịp tim vẫn còn đập loạn trong lồng ngực.

``\dots Đã đủ\dots\ xa chưa?'' tôi hỏi, giọng vẫn còn vương chút run rẩy.

``Anh nghĩ là\dots\ đủ rồi,'' Nii-chan đáp thay cho cả ba, mắt dõi từ phía chúng tôi vừa tới, như để chắc rằng không có ai đuổi theo.

Tôi siết chặt bàn tay mình. Trong lòng tôi chỉ còn lại một mớ hỗn độn: cái quái gì vừa xảy ra thế? Tại sao Saya-san, Kaoru-san và Miku-san lại\dots\ lặp lại từng câu từng chữ, y như những con rối vô hồn?

Và điều khiến tôi lo lắng hơn cả, là cách họ phản ứng mỗi khi chúng tôi nhắc đến Akane-san và Yuka-san. Họ như đang cố tình né tránh, lảng sang chuyện khác, hoặc trả lời một cách chung chung, như thể không muốn chúng tôi đào sâu thêm về hai người bạn đó. Tôi không thể không nghĩ rằng\dots\ có điều gì đó không ổn đã và đang xảy ra với hai người ấy.

``Onee-sama.'' tiếng Reiko cất lên nhẹ nhàng, kéo tôi khỏi dòng suy nghĩ.

Tôi quay sang, thấy em ấy lục trong túi đeo, lấy ra một chai nước và đưa cho tôi. ``Chị uống chút nước đi.''

``Cảm ơn em, Reiko.''

Tôi khẽ gật đầu nhận lấy, uống một ngụm để làm dịu cổ họng khô khốc.

Ngay sau đó, em ất lại quay sang Nii-chan: ``Onii-sama.'' Lần này đưa thêm một chai nước khác.

Anh hơi khựng lại, rõ ràng chưa quen với cách xưng hô này, nhưng cuối cùng vẫn cầm lấy. ``\dots Anh xin.''

Tôi mỉm cười, khẽ khen: ``Em chu đáo thật đấy, Reiko.''

``Ehehe\~{}\dots'' Reiko chỉ cười, ôm sát chiếc túi của mình, rồi cũng lấy ra một chai khác cho riêng mình.

Nhìn nụ cười trong trẻo ấy, tôi chợt thấy lòng ấm lại. Dù cho thế giới xung quanh có điên đảo thế nào, thì tình cảm gia đình vẫn luôn là điểm tựa vững chắc nhất. Như lúc này đây, chỉ một cử chỉ quan tâm đơn giản của Reiko cũng đủ xua tan đi phần nào cảm giác bất an trong lòng tôi.

Tôi khẽ xoa đầu em gái, cảm nhận được hơi ấm quen thuộc truyền qua những ngón tay. Trong khoảnh khắc ấy, ranh giới giữa thực và mơ dường như không còn quan trọng nữa. Chỉ cần có nhau, chúng tôi sẽ vượt qua được tất cả.

Nhưng rồi, chúng tôi vẫn phải đối mặt với thực tại đầy tàn khốc. Những sự ``bình thường'' ấy hóa ra cũng rất xa xỉ với tình hình hiện tại đó, khi bủa vây xung quanh chúng tôi chỉ có đầy rẫy sự ``bất thường.''

Nii-chan lên tiếng, giọng trầm ngâm: ``Saikyou, Reiko\dots\ giấc mơ ngôi trường sụp đổ mà chúng ta thấy trước đây\dots\ anh nghĩ nó không phải chuyện vô nghĩa. Liên kết với chuyện vừa rồi, có khi nào\dots\ giấc mơ đó đã tràn sang hiện thực?''

Tôi cứng người lại. Ý nghĩ ấy nghe thật điên rồ\dots\ nhưng với những gì chúng tôi đã trải qua, tôi không thể phản bác.

Reiko ôm chặt chai nước, giọng hơi run: ``Nhưng vấn đề là\dots\ Akane-san và Yuka-san rốt cuộc ở đâu? Chúng ta còn chẳng biết gì nhiều về họ, ngoài tên và mối quan hệ thân thiết của họ nữa.''

``Đó chính là vấn đề đấy.'' Nii-chan thở dài, chống tay lên trán. ``Ta chẳng thể nào biết được manh mối nếu không hiểu rõ tính cách của họ đâu.''

Tôi im lặng một lúc, rồi nói chậm rãi: ``\dots Chỉ còn cách đi một vòng quanh thị trấn thôi. Biết đâu có người dân ở đây từng gặp họ.''

Cả hai im lặng nhìn tôi, rồi cùng gật đầu. Không còn lựa chọn nào khác.

``Vậy thì\dots'' tôi hít sâu, cố lấy lại chút sức lực, chuyển bị cho cuộc hành trình dài hơi, ``chúng ta bắt đầu tìm thôi.''

\ct

Chúng tôi quyết định đi dọc theo con phố truyền thống vừa rời khỏi ban nãy. Những căn nhà gỗ thấp tầng, những tấm biển cũ kỹ treo lủng lẳng trước cửa tiệm, tất cả vẫn y nguyên như trước, nhưng Saya-san, Miku-san và Kaoru-san thì không còn ở đó nữa, có lẽ vừa mới rời.

Thế nhưng, bầu không khí vắng lặng ấy khiến lưng tôi bất giác rợn lạnh. Không hiểu tại sao, nhưng tôi có cảm giác như ba người họ \emph{chưa từng tồn tại}, chỉ để lại một dư chấn vô hình trong ký ức chúng tôi.

*SOẠT SOẠT* *TÁCH!*

Đang định bước đi, tôi chợt thấy một cặp tai thỏ lấp ló từ xa. Tôi dừng lại, nheo mắt nhìn kỹ hơn. Cặp tai ấy lổn nhổn một lúc, rồi từ từ nhô lên cao. Đó là một cô gái có mái tóc xanh dương nhạt, buộc thành hai bím đuôi sam, mặc bộ đồng phục giống hệt chúng tôi, đang cầm máy ảnh chụp ở một góc phố. Cô ấy di chuyển nhẹ nhàng, uyển chuyển giữa những gian hàng, thỉnh thoảng dừng lại để ngắm nghía một góc nào đó qua ống kính.

Có gì đó rất khác biệt ở cô gái ấy. Không giống như vẻ trống rỗng của ba cô bạn mà chúng tôi gặp, ánh mắt cô ấy sáng lên đầy sức sống. Nụ cười tươi tắn khi bấm máy, dáng người năng động khi di chuyển - tất cả đều toát lên một sự tươi tỉnh đến lạ kỳ giữa bầu không khí nặng nề này.

``Saikyou!'' giọng Nii-chan vang lên từ phía sau, kéo tôi về với thực tại. ``Đừng đứng đó nữa, chúng ta còn phải tìm manh mối về Yuka-san nữa mà.''

``V-Vâng!'' tôi đáp, vội vàng rời đi.

\ct

Chúng tôi ghé vào một quán cà phê nhỏ đầu phố, bắt đầu cuộc tra hỏi. Khi Nii-chan lịch sự hỏi thăm về Akane-san và Yuka-san, một người đàn ông trung niên chỉ ngẩng đầu khỏi tờ báo, cau mày: ``Akane? À\dots\ cô gái tóc nâu sáng, hình như tôi có thấy qua vài lần. Nhưng Yuka? Nó là ai vậy?''

Ông ta nghiêng đầu, mắt lộ rõ vẻ khó hiểu. Tôi tưởng rằng việc hai người ấy đi cùng với nhau thì sẽ khiến cho ông ta nhận ra, nhưng không, tôi đã lầm.

Mà ngẫm lại, Akane-san có bao giờ đi một mình không nhỉ?

``Dạ, cảm ơn bác ạ,'' anh trai tôi khẽ gật đầu đáp, rồi hướng ra về phía chúng tôi, lắc đầu.

Chúng tôi sang bàn bên, hỏi một nhóm phụ nữ đang tán gẫu. Một người nhíu mày, đáp nhanh: ``Chưa từng nghe tới. Ở thị trấn này làm gì có người tên như thế.''

Giọng điệu dứt khoát, như thể câu hỏi của chúng tôi hoàn toàn vô nghĩa.

Reiko cắn môi, cố gắng lấy hết can đảm hỏi thêm: ``Thế còn Akane-san\dots\ chị có biết gì về cô ấy không?''

Một người khác khẽ cười gượng: ``À, cô bé ấy\dots\ cũng có lần thấy, đi ngang qua thôi. Nhưng còn bạn đi cùng thì\dots\ chị cũng không chắc. Hình như cô ấy chỉ đi một mình mà?''

Tôi thoáng ngỡ ngàng, quay sang nhìn Nii-chan. Ánh mắt anh cũng hiện rõ sự cảnh giác. Không lẽ\dots\ tất cả bọn họ đều thật sự không biết Yuka-san sao?

Chúng tôi thử sang một cửa hàng tiện lợi gần đó. Cậu nhân viên trẻ tuổi đang sắp xếp kệ hàng, nghe xong câu hỏi thì bối rối gãi đầu: ``Xin lỗi\dots\ mình không biết ai như vậy. Khách quen của chúng mình cũng không có tên đó đâu.''

``Vậy\dots\ cậu đã từng gặp vị khác nào có bờm mèo, tóc tím ngắn không?'' tôi hỏi, toát mồ hôi. Người nhân viên kia lắc đầu.

Thái độ ấy không giống dối trá. Nó giống như\dots\ thật sự chẳng hề tồn tại người tên Yuka trong ký ức của cậu ta.

Điều đó dấy lên trong lòng tôi một câu hỏi: ``Liệu cái tên Nekomiya Yuka đó có thực sự sống ở thị trấn này không?''

Nhưng rồi, tôi gạt ý nghĩ đó sang một bên, tiếp tục cùng hai người lùng sục khắp thị trấn để tiếp tục tra hỏi mọi người, với niềm tin rằng ít nhất ai trong số họ đều biết cái tên ``Yuka'' đó.

\ct

Một tiếng đồng hồ trôi qua. Chúng tôi đi từ quán này sang quán khác, từ người bán hàng rong ngoài phố đến ông chủ cửa tiệm đồ lưu niệm. Cùng một câu hỏi, cùng với từng chi tiết mà chúng tôi biết về cô bạn tai mèo kia.

Nhưng, câu trả lời mà chúng tôi nhận về vẫn chỉ xoay quanh ba kiểu: chưa từng nghe, có gặp Akane nhưng chỉ một mình, hoặc nghiêng đầu khó hiểu trước cái tên Yuka. Thậm chí, một số người còn bảo chúng tôi bị điên.

Sau mỗi lần như vậy, bầu không khí giữa ba chúng tôi lại nặng nề thêm một chút. Tôi thấy bàn tay Reiko nắm chặt gấu váy đến trắng bệch, đôi mắt thì ỉu xìu hiện rõ sự thất vọng pha lẫn sự hoang mang. Còn Nii-chan, vẻ mặt bình tĩnh thường ngày giờ đây cũng phảng phất nỗi lo âu, mắt khẽ nheo lại như đang ghép nối những mảnh ghép không vừa khít.

Câu hỏi ban nãy về sự tòn tại của cô ấy bỗng chốc lại hiện về. Nỗi hoang mang ấy, một khi đã chạm tới, không cách nào gạt bỏ được.

\pov{Haragen Reiko}

Bước chân tôi dần nặng nề sau hơn một giờ dò hỏi vô vọng. Bên cạnh, Onee-sama và Onii-sama vẫn cố giữ vẻ cứng cỏi, nhưng tôi biết trong lòng cả hai cũng đang hoang mang chẳng kém gì mình.

Từ sáng tới giờ, mọi thứ cứ như một cơn ác mộng vậy. Giấc mơ kỳ lạ của Onee-sama về con tàu bí ẩn, rồi đến việc ba người bạn mới quen của chúng tôi đột nhiên quên hẳn Yuka-san. Tất cả những điều đó khiến tôi cảm thấy bất an đến lạ. Tôi vẫn nhớ rõ nụ cười ấm áp của Yuka-san hôm qua, vậy mà giờ đây, dường như cô ấy chưa từng tồn tại trong ký ức của mọi người. Cảm giác như thể có một bàn tay vô hình nào đó đang xóa đi từng mảnh ký ức, từng dấu vết của sự hiện diện.

Chính vào lúc ấy, bụng tôi bất chợt réo lên một tiếng rõ to, đến mức cả hai chị ngoảnh về phía tôi. Tôi giật nảy, mặt nóng bừng vì ngượng.

``\dots Xin lỗi, Onee-sama, Onii-sama\dots'' tôi cúi gằm mặt, vô thức che bụng lại.

Cả hai thoáng khựng, rồi khẽ bật cười. Onee-sama xoa nhẹ đầu tôi, giọng dịu dàng: ``Không sao đâu, Reiko. Chắc là chúng ta cũng nên nghỉ chân một chút nhỉ? Chúng ta đi bộ cũng lâu rồi mà.''

``Có một quán ăn ở ngay trước mặt kìa,'' anh tôi đi ở phía trước, đôi mắt hướng về quán ăn cách chúng tôi chỉ vài trăm mét. Một quán ăn nhỏ, không bảng hiệu cầu kỳ, chỉ có vài bộ bàn ghế đơn giản bên trong. Hơi ấm và mùi hương đồ ăn tỏa ra khiến tôi thấy dễ chịu, như thể đã được an ủi sau cả buổi sáng nặng trĩu.

Chúng tôi ngồi xuống một bàn gần cửa sổ, thở phào nhẹ nhõm sau quãng đường đi bộ. Một bà chủ quán tầm trung niên tiến lại gần, nụ cười hiền hậu trên môi.

``Các cháu muốn gọi món gì?'' bà ấy hỏi, tay cầm sẵn cuốn sổ nhỏ.

``Dạ cho cháu một suất mì ramen,'' Onee-sama nói.

``Cháu cũng vậy ạ,'' Onii-sama gật đầu.

Tôi nhìn qua thực đơn, rồi ngẩng lên: ``Cho cháu một suất mì ramen và một suất cơm cà ri ạ!''

``Hả!?'' Onii-sama giật mình, quay sang nhìn tôi với ánh mắt ngạc nhiên. Onee-sama thì bật cười khúc khích sau tay.

Bà chủ quán cũng tròn mắt: ``Ôi trời, cháu bé nhỏ nhắn thế này mà ăn được nhiều vậy sao?''

Tôi gãi má, hơi ngượng: ``Dạ\dots\ vì\dots\ cháu đói quá ạ.'' Thực chất, tôi là một người phàm ăn, nhưng tôi biết một người con gái sẽ không bao giờ thừa nhận điều đó với những người ngoài. Nhất là khi Onii-sama cũng ở đây nữa chứ!

Trong lúc ngồi chờ món, đầu óc tôi bất giác trôi ngược về ngày hôm qua. Từ khoảnh khắc chúng tôi ở trên con tàu lạ tới Aogami, rồi giấc mơ của chị, cho tới khi vướng vào vòng lặp cùng những gương mặt trống rỗng của ba cô bạn cùng lóp với họ. Tôi siết chặt hai tay lên đùi, nỗi bất an dâng cao.

Nhưng trong dòng suy nghĩ hỗn độn ấy, một điểm sáng bỗng lóe lên. Ngôi đền Aogami. Nơi mà Saya-san và Akane-san đã nhắc tới - dù tôi không rõ là do tự họ nhắc tới, hay do một thế lực nào đó ép họ nói vậy. Đấy là nơi mà chúng tôi đã thấy Sora-san cầu nguyện một cách khác thường. Và cũng là nơi Shino-san đứng trò chuyện với chúng tôi trước khi dịch chuyển tới đây. Tất cả những sự kiện đó dường như đều có liên quan tới nơi thiêng ấy.

Có lẽ\dots\ đó có thể là manh mối để giải đáp những điều bí ẩn đang xảy ra. Hoặc chí ít, là một thông tin nào đó khác hữu ích.

``\dots Onee-sama, Onii-sama.'' Tôi ngập ngừng cất tiếng.

``Hửm?'' Hai người đồng loạt nhìn sang.

``Em nghĩ\dots\ chúng ta nên thử đến đền Aogami.''

Ngay lập tức, nét mặt của hai anh chị nhăn lại. Không cần tôi giải thích gì thêm, họ cũng ngầm hiểu được ý định của tôi là gì.

``Em có chắc là nơi đó sẽ cho ta điều gì có ích chứ?'' Onii-sama hỏi, giọng có chút ngờ vực.

``Không hẳn ạ\dots'' tôi lắc đầu, nhưng rồi giương ánh mắt kiên định, ``Em chỉ biết là\dots\ linh cảm mách bảo rằng ta cần đến đó.''

Onii-sama nhíu mày, suy tư. Trong ánh mắt anh thoáng hiện lên sự lo lắng xen lẫn với quyết tâm, như thể anh cũng đã nghĩ tới điều này từ lâu.

Onee-sama cũng gật nhẹ, đôi mắt chị tối lại khi những ký ức về giấc mơ kỳ lạ ban sáng ùa về. Không khí bàn ăn trầm xuống, nặng trĩu bởi những suy tư chưa thể nói thành lời, nhưng tôi cảm thấy ít ra mình đã không ngồi im lặng mãi.

Đôi mắt tôi chợt để ý trên kệ có xếp mấy tờ báo cũ, giấy đã ngả vàng. Một dòng tít lớn lọt vào mắt tôi:

\begin{center}
  \uline{Hàng loạt vụ mất tích bí ẩn gần nhà ga Aogami.}
\end{center}

Tôi run tay cầm lên, đọc thoáng qua. Nội dung ngắn gọn nhưng đầy ám ảnh: trong nhiều tháng qua, số người mất tích tăng đột biến, tất cả đều có điểm chung: biến mất gần khu vực nhà ga Aogami. Một vài nhân chứng kể rằng, vào đêm muộn họ nghe thấy tiếng còi tàu mà họ chưa từng nghe bao giờ, nhưng cảnh sát chưa xác thực được điều gì.

Tôi đem tờ báo ra quầy, hướng về bà chủ quán, người vẫn đang làm món cho khách.

``Bác ơi\dots\ cháu đọc thấy vụ này\dots\ có thật không ạ? Người dân thật sự nghe thấy\dots\ tiếng còi tàu sao?''

Bà ấy ngước về phía tôi, trầm ngâm một thoáng, rồi hạ giọng, mắt tiếp tục nhìn vào bát mì đang làm dở: ``À\dots\ chuyện ấy à. Có vài người kể lại, bảo là nửa đêm nghe thấy tiếng còi tàu dài, nằm vất vưởng đâu đó ở gần nhà ga Aogami. Chính tai bác cũng đã nghe thấy nó rồi.''

``Tiếng còi ạ?''

``Nó không giống mấy loại tàu điện bây giờ đâu,'' bà đáp, đặt bát mì lên khay. ``Nó nghe như tiếng tàu chạy than củi, hồi mấy chục năm trước. Bác thậm chí còn nhìn thấy khói trắng phun ra từ trong nhà ga ấy, điều mà bác đã không thấy nó từ hàng chục năm rồi. Khung cảnh ấy\dots\ một khi đã nghe rồi thì không quên được.''

Tôi thấy tim mình lạnh toát. Một hồi ức mơ hồ vụt hiện: khói trắng, bánh sắt nghiến trên đường ray, và bóng lưng Yuka-san đi vào khoang tàu\dots

Không khí trong quán chùng xuống, chỉ còn tiếng đũa thìa leng keng.

Khi đồ ăn được bưng ra, chúng tôi cố gắng ăn trong yên lặng, dù chẳng ai nuốt nổi, khi trong đầu chúng tôi vẫn còn quanh quẩn những bí ẩn xung quanh nhà ga ấy.

\ct

Mười hai giờ trưa.

Dứt bữa, Onii-sama đặt đôi đũa, ánh mắt dứt khoát: ``Đi thôi. Tới đền Aogami. Nếu đó là nơi gắn liền với tất cả manh mối, chúng ta không thể bỏ qua.''

``Em cũng nghĩ vậy,'' tôi khẽ gật đầu, giọng trầm xuống. ``Những gì chúng ta vừa chứng kiến\dots\ không thể để mặc được.''

Onee-sama siết chặt nắm tay: ``Phải tìm ra sự thật. Về Yuka-san, về những người mất tích, và cả về ngôi đền đó nữa.''

``Chúng ta sẽ làm được thôi,'' tôi mỉm cười, cố trấn an chị mình dù trong lòng cũng đang bất an không kém.

Cả ba bước ra khỏi quán, ánh nắng trưa đã chói chang hơn. Để tiết kiệm thời gian và công sức, chúng tôi quyết định bắt một chuyến xe buýt hướng về ngôi đền.

Bánh xe lăn chầm chậm trên con đường gập ghềnh, mà trong lòng tôi chỉ có một câu hỏi dằn vặt:

``Liệu tiếng còi tàu kia\dots\ có đang đợi chúng tôi ở phía trước?''

\ct

Một giờ chiều. Chiếc xe buýt vừa rời đi, chúng tôi bước xuống trước một con dốc thoai thoải. Trước mắt trải ra từng bậc đá cũ kỹ, sứt mẻ theo năm tháng, dẫn thẳng lên núi. Tôi ngẩng nhìn, cơn gió cuối đông rít qua khe áo len, kéo theo mùi ẩm ướt của sương mù đang vương trên cao.

``\dots Chính là đây sao?'' tôi lẩm bẩm, đôi mắt dõi theo những bậc đá uốn khúc.

Một cánh cổng torii đỏ sẫm hiện ra ở đầu dốc, vững chãi mà trầm mặc. Và rồi, giống như một dải màu rực rỡ chen giữa nền xám xịt của trời đất, những bụi hoa bỉ ngạn đỏ mọc san sát hai bên đường. Cánh hoa đỏ như máu, rũ xuống từng chùm, tựa những ngọn lửa âm ỉ cháy dưới lớp tro tàn.

Tôi khẽ siết chặt dây quai túi. Đây không phải lần đầu tiên chúng tôi bắt gặp chúng. Những cánh hoa ấy từng vây quanh chúng tôi trong cơn ác mộng, lẫn trong khung cảnh sụp đổ, trong giấc mơ lặp đi lặp lại, và cả hôm ấy, khi cùng Sora-san thoát khỏi trận sạt lở. Lần nào cũng vậy, sự xuất hiện của chúng như một lời cảnh báo, một điềm gở chẳng lành.

Onee-sama khựng lại một nhịp, nhìn chằm chằm vào một khóm hoa ngay sát chân mình. ``Chúng\dots\ vẫn ở đây.'' giọng chị khẽ run.

Onii-sama chép miệng, cố gắng tỏ ra bình thản: ``Ừ, nhưng lần này ít nhất trời còn sáng, và anh vẫn có thể đi được.. Lần trước\dots\ thì anh còn phải để Saikyou cõng nữa kìa.''

Tôi tròn mắt quay sang, và đúng như dự đoán, Onee-sama nhoẻn cười đầy ẩn ý: ``À phải rồi, Nii-chan. Lúc đó anh chẳng còn cãi lại được em nữa cơ.''

Khuôn mặt Onii-sama thoáng đỏ lên, cậu quay đi, khẽ hắng giọng: ``Ờ thì\dots\ hôm đó anh chấn thương nặng, chứ không thì\dots''

Cả tôi và Onee-sama cùng phì cười. Giữa bầu không khí âm u, sự đùa vui nhỏ ấy khiến lòng tôi dịu lại đôi chút.

Chúng tôi tiếp tục bước đi, bậc đá lạo xạo dưới chân, xen lẫn tiếng gió thổi hun hút từ rừng thông ven sườn núi. Vài du khách lác đác đi ngược chiều, ánh mắt vô hồn, lặng lẽ bước như những chiếc bóng không cảm xúc. Tôi thoáng rùng mình, họ giống hệt những nhân vật ngoài lề trong trò chơi, chỉ hiện hữu để lấp đầy khoảng trống, chẳng màng đến sự tồn tại của chúng tôi.

``Đến cả những người ở đây cũng hầu như không có cảm xúc gì luôn\dots'' tôi lẩm bẩm, cảm giác như thể cả thế giới này chỉ có chúng tôi thực sự sống vậy.

Tôi ngẩng nhìn con đường trải dài dưới sương mù mỏng, những cánh bỉ ngạn đỏ cứ nối tiếp nhau bất tận. Mỗi bước tiến lên, trái tim tôi lại dội một nhịp dồn dập, như thể chính con đường này đang níu kéo ký ức méo mó kia trở về.

Và rồi, sau khi vượt qua thêm một cánh cổng torii, bóng dáng mái ngói cong cong của ngôi đền Aogami hiện dần ra, mờ ảo trong màn sương lạnh.

Không gì khác hơn, đó là ngôi đền mà Sora-san đã tới để cầu nguyện trong nước mắt.

Nhưng giờ, chúng tôi tới đây với một mục đích khác: tìm xem có manh mối nào ở đây không.

Chúng tôi chia nhau đi từng ngóc ngách của ngôi đền.

``Em sẽ kiểm tra hành lang bên trái,'' tôi nói, nhìn về phía lối đi tối om.

Onee-sama gật đầu: ``Chị sẽ vòng ra phía sau hậu điện.''

``Còn anh sẽ xem xét khu vực thờ chính,'' Onii-sama nói ngắn gọn, bước về phía trước.

Tôi lần theo hành lang, mỗi bước đi đều cẩn thận. Tay tôi lướt trên từng bề mặt gỗ, từng tấm bảng cũ kỹ. Không có gì đặc biệt ngoài những vết xước thời gian và những sợi dây buộc đã bạc màu.

``Có phát hiện gì không?'' Onee-sama gọi từ xa.

``Chưa ạ,'' tôi đáp lại. ``Em chỉ thấy toàn vết tích cũ thôi.''

``Bên này cũng vậy,'' giọng anh vọng lại từ gian thờ chính.

Khoảng hơn nửa giờ sau, khi ánh nắng trưa đã dịu bớt, tôi quyết định quay lại gian thờ chính để nghỉ ngơi. Onee-sama và Onii-sama cũng đã ở đó.

``Em mệt rồi à?'' chị tôi hỏi nhẹ nhàng.

Tôi gật đầu, định trả lời thì bỗng mắt tôi bắt gặp một điều gì đó. Phía sau pho tượng thần trên bệ thờ, có một khe hở nhỏ mà trước đó tôi chưa từng để ý.

``Có cái gì đó ở đằng kia\dots'' tôi nói thầm, nhưng chừng đó vừa đủ đề hai anh chị có thể nghe thấy.

``Reiko?''

Linh cảm trong lòng bất chợt thôi thúc, tôi cúi xuống, thò tay vào thử. Ngón tay chạm vào một thứ cứng, mát lạnh. Tôi kéo nó ra, trên tay là một khối hình chữ nhật nhỏ, không lớn hơn lòng bàn tay. Trên bề mặt có dán nhãn cảnh báo tái chế, kèm theo một mã QR in mờ.

``Onii-chan, Onee-chan! Lại đây xem cái này!'' tôi gọi.

Cả hai nhanh chóng bước tới. Onii-sama nheo mắt nhìn kỹ vật thể một lúc, rồi khẽ lẩm bẩm: ``Trông nó giống\dots\ một cục pin.''

Ngay khi câu nói ấy vang lên, tôi thấy vẻ mặt cả hai cứng lại. Như thể có điều gì đó lóe lên trong ký ức của họ.

Tôi không hiểu, chỉ đứng lặng nhìn.

``Onee-sama, chị đang làm gì vậy?'' tôi hỏi khi thấy chị vội vã lục tìm gì đó trong túi áo.

``Đợi chị một chút,'' Onee-sama nói, giọng có chút gấp gáp. ``Chị nghĩ mình đã hiểu ra điều gì đó.''

Chị lấy ra một mảnh bo mạch chủ - thứ mà chị từng nhặt ở căn phòng đầy sắc đỏ trong vòng lặp trước. Chị đặt khối pin vào vị trí trống đó một cách đầy tỉ mỉ. Vừa khít, như thể vốn dĩ được chế tạo để gắn với nhau.

``Nii-chan,'' Onee-sama quay sang anh trai song sinh của mình, ``anh có nghĩ như em không?''

Onii-sama gật đầu, ánh mắt sáng lên: ``Ừ. Nó giống hệt bộ phận trong chiếc đồng hồ đeo tay đó.''

``Không thể sai được\dots'' Onee-sama thì thầm.

Tôi đứng giữa hai người, cố gắng hiểu những gì đang diễn ra. Họ dường như quá để ý tới cái bộ phận nhìn bên ngoài có vẻ vô tri ấy, còn trong khi đó tôi là kẻ đứng ngoài--

\dots

\dots Chờ đã.

Bất chợt, một mảnh ký ức vụt qua trong đầu: cũng ở chính ngôi đền này, nhưng ở một không gian khác, khi tôi đoàn tụ - dù rằng cậu ấy đã chết - đã từng nhắc chúng tôi về ``chiếc đồng hồ đeo tay''. Cậu ấy cũng nhắc tới việc nó rất quan trọng.

``Em nhớ rồi!'' tôi buột miệng. ``Shino-san đã nói chiếc đồng hồ là chìa khóa để giải đáp tất cả, phải không ạ?''

``Đúng vậy,'' Onii-sama trầm ngâm. ``Về sự tồn tại của chúng ta, về việc tại sao chúng ta lại ở đây\dots''

Lúc đó, tôi không hiểu, và cả bây giờ cũng vẫn còn mơ hồ, nhưng sự trùng hợp khiến lòng tôi bất giác thắt lại. Có lẽ những gì mà cậu ấy nói hôm đó\dots\ không phải là vô nghĩa.

``Chúng ta nên giữ cẩn thận cục pin này,'' Onee-sama quyết định. ``Nó có thể là một manh mối quan trọng.''

``Vâng,'' tôi gật đầu. ``Nhưng chúng ta vẫn nên tiếp tục tìm kiếm thêm manh mối khác.''

Onii-sama đồng ý: ``Được rồi. Hãy chia nhau ra tìm kiếm thêm.''

Rồi, cả ba chúng tôi lại tiếp tục tìm kiếm, dạo quanh khuôn viên của ngôi đền, và thậm chí soi xét từng bụi rậm quanh nó.

Suốt cả buổi chiều, ba chúng tôi đi khắp từ trước ra sau, quanh quanh sân đền, thậm chí cả nhà kho nhỏ ở góc - tôi dám chắc nó đã bị bỏ hoang bị lâu, vì khi chúng tôi vào, thứ chúng tôi thấy chỉ là một đống bụi bặm và đầy rẫy mạng nhện - nhưng chẳng còn gì khác.

``Chẳng có gì cả\dots'' Onee-sama thở dài, phủi bụi bám trên tay áo. ``Ngoài cái cục pin đó ra thì\dots\ nơi này không có một cái gì đó có thể giúp chúng ta cả.''

``Có lẽ chúng ta nên nghỉ thôi,'' Onii-sama lên tiếng, ánh mắt vẫn dõi theo những tia nắng đang dần hạ xuống len qua khe cửa nhà kho. ``Trời cũng xế rồi.''

Tôi khẽ gật đầu, dù trong lòng vẫn còn chút tiếc nuối: ``Em cũng nghĩ vậy. Nhưng mà\dots\ cái cục pin đó\dots''

``Ừm,'' Onee-sama nắm chặt nó trong tay, ``chị cũng thấy nó rất đáng ngờ. Một thứ như vậy\dots\ lạitrong một ngôi đền cổ\dots''

``\dots Quả thật là rất kỳ lạ,'' tôi nói. Cả hai anh chị đều có cùng quan điểm với tôi.

Tới gần bốn giờ, khi ánh nắng đã bắt đầu ngả dần về phía tây, chúng tôi đành dừng lại. Chúng tôi bước ra khỏi gian chính, ánh nắng cuối buổi chiều đã hắt xiên qua rặng cây phía xa. Tôi khẽ thở dài, định cùng hai người rời khỏi ngôi đền thì bỗng gió lặng hẳn. Bầu không khí nặng xuống, mọi âm thanh như biến mất.

Ngay trước mặt chúng tôi, nơi bậc thềm đá, một dáng người từ từ hiện ra. Mái tóc nâu dài tới ngang lưng, cái cài tóc kia, đôi mắt xanh tĩnh lặng, cùng bộ đồng phục màu xanh đậm ấy\dots

Tôi nhận ra ngay, là Shino-san.

Cô ấy đứng đó, đôi mắt nhìn thẳng chúng tôi. ``Các cậu\dots\ tình hình thế nào rồi?'' giọng cô vang lên, nhẹ mà rõ ràng.

Onii-sama giọng trầm xuống, thở dài một tiếng tỏ rõ sự mệt mỏi: ``Chúng tôi vẫn chưa tìm thấy Sora. Nhưng\dots''

Anh lấy từ trong túi ra cục pin, giơ lên cho Shino thấy. ``Đó là một một phần của đồng hồ đeo tay, chúng tôi thấy ở đằng sau bức tượng kia.''

``Chúng mình tìm thấy thứ này ở gần nhà ga,'' tôi tiếp lời. ``Mình không chắc nó có ý nghĩa gì không, nhưng\dots\ theo những gì cậu nói trong vòng lặp trước\dots\ nó có thể giúp ích cho ba chúng mình.''

``Đúng là mình nói như vậy,'' Shino-san nói. ``Nhưng ngoài cái đó ra, các cậu còn trải qua những gì khác?''

Onee-sama lập tức chen vào: ``Shino-san, đêm qua, tôi đã thấy một giấc mơ. Tôi\dots\ không chắc có nên gọi nó là giấc mơ nữa.''

Đôi mắt xanh ấy khẽ lay động. ``Cậu mơ thấy gì?''

Thế là cả ba chúng tôi thay phiên kể lại, không bỏ sót chi tiết. Đặc biệt là về con tàu hơi số hiệu 211493, về Yuka-san bước vào con tàu, biểu hiện trông rất lạ thường như vừa bị thôi, và rồi biến mất. Shino-san lắng nghe chăm chú, ánh mắt dần trở nên u ám.

Tôi bổ sung thêm: ``Và còn nữa\dots\ sáng nay, khi ba bọn mình tới điểm mà Akane-san hẹn, Kaoru-san, Saya-san, Miku-san\dots\ họ không còn nhớ gì về Yuka-san nữa. Như thể cô ấy chưa từng tồn tại vậy.''

Cô bạn cũ im lặng một lúc lâu. Cô nhìn cục pin trong tay, rồi ngẩng lên, giọng trầm xuống: ``Đó không phải giấc mơ. Đó là ký ức đang bị bóp méo bởi vòng lặp thời gian. Quá khứ năm 1946 và hiện tại của các cậu đang chồng lấn lên nhau, tạo ra những khoảng trống trong ký ức mọi người.''

Cô dừng lại một chút, rồi tiếp tục: ``Đây là một phép thử\dots\ để xác định sức mạnh của ước nguyện mà Sora để lại. Và có vẻ như\dots\ nó đã bắt đầu tác động đến thực tại rồi.''

Tôi giật mình. Anh chị lớn của tôi cũng không giấu nổi vẻ bối rối.

Shino-san tiếp lời: ``Đó cũng là lý do vì sao mọi người trong thị trấn chẳng ai nhớ Yuka. Trong vòng lặp, ký ức của họ bị viết lại, xóa bỏ sự tồn tại của cô ấy.''

``Giống như một cuốn sách bị xé đi những trang quan trọng,'' cô ấy nói tiếp, giọng trầm xuống. ``Những ký ức về cô ấy đã bị xóa sạch, thay thế bằng những mảnh ghép giả tạo. Và điều đáng sợ là\dots\ họ thậm chí không nhận ra điều đó.''

Lời giải thích của Shino-san khiến cả ba chúng tôi đều sững người. Onii-sama khẽ nhíu mày, ánh mắt lộ rõ vẻ căng thẳng khi anh cố gắng tiếp nhận thông tin này. Onee-sama thì siết chặt nắm tay, môi mím lại thành một đường thẳng, như thể chị đang cố kìm nén cảm xúc trước sự thật đáng sợ này. Còn tôi, tôi cảm thấy một cơn ớn lạnh chạy dọc sống lưng - ý nghĩ về việc ký ức của mọi người có thể bị xóa sạch và thay thế như vậy khiến tôi không khỏi rùng mình.

``Nhưng vẫn có ngoại lệ\dots'' cô nói, hướng mắt về phía chúng tôi. ``Những người như ba cậu, những kẻ không bị cuốn hoàn toàn vào dòng chảy giả tạo, vẫn còn giữ được mảnh ghép thật sự.''

Tim tôi đập nhanh hơn. Vậy ra chúng tôi là ngoại lệ, một trong số ít những kẻ còn nhớ.

Shino-san nhìn tôi, rồi quay sang hai anh chị: ``Nếu muốn tìm đường thoát khỏi vòng xoáy này, các cậu phải lần theo dấu vết của Yuka và Akane. Hãy tìm những người từng thật sự gần gũi với họ. Chỉ vậy, ký ức bị chôn vùi mới có thể gợi mở.''

Tôi ngập ngừng, rồi buột miệng: ``Nhưng Kaoru-san, Miku-san, và Saya-san đều từng là bạn học của họ. Vậy mà\dots\ chẳng ai nhớ gì về Yuka-san hết. Nếu vậy thì chúng mình phải bắt đầu từ đâu?''

Cậu ấy khẽ nhắm mắt lại, nét mặt thoáng một nụ cười buồn: ``Những ký ức bị xóa đi có thể không trở lại trong cùng một hình dạng. Đừng chỉ hỏi họ về Yuka như một cái tên. Hãy tìm những mảnh ghép khác. Ký ức có thể ẩn dưới dạng một chi tiết nhỏ, nhưng vẫn tồn tại. Nơi nào còn dấu vết, nơi đó vẫn còn con đường.''

``Và còn một điều nữa,'' cô ấy nói thêm, giọng trở nên nghiêm trọng. ``Tối nay, các cậu hãy đến ga Aogami. Đúng lúc khi ngoài đường đã vãn dần. Có thể các cậu sẽ tìm thấy thứ gì đó\dots\ hoặc ai đó.''

Chúng tôi chưa kịp hỏi thêm thì thân ảnh Shino bắt đầu nhạt dần, như sương tan vào gió. Đôi mắt xanh nhìn chúng tôi lần cuối, rồi biến mất hẳn.

Chỉ còn lại ba người đứng giữa sân đền vắng, mang theo manh mối về bộ phận của chiếc đồng hồ và những lời giải thích của cô ấy về mối liên hệ giữa nó với những sự kiện đang diễn ra.

Chúng tôi rời khỏi sân đền, bóng nắng đã ngả dần về phía tây. Con đường bậc đá phủ sương buổi chiều như kéo dài bất tận, từng bước đi nghe khẽ vang vọng trong không khí tĩnh lặng.

Onii-sama lên tiếng trước, giọng trầm xuống: ``Shino-san nói rõ rồi. Nếu muốn hiểu chuyện này, chúng ta phải đến ga Aogami tối nay.''

Tôi nhìn thoáng qua anh tôi, thấy khóe môi anh mím chặt. Onee-sama bước chậm lại, tay khoanh trước ngực: ``Đành phải vậy chứ làm sao. Em cảm thấy rằng cô ấy không nói dối đâu..''

Tôi khẽ gật đầu. ``Em cũng nghĩ vậy. Ít nhất, thử xem có gì lạ không.''

Anh thở ra, một hơi dài. ``Được, vậy tối nay. Chúng ta ăn tối, rồi quay lại ga. Tầm khoảng chín, mười giờ đêm là ổn nhất.''

Tôi chỉ ``dạ'' một tiếng. Theo như tôi biết, buổi đêm luôn là khoảng thời gian mà những hiện tượng ma mị hoạt động mạnh nhất, nhưng khoảng giờ đó cũng đủ để không bị coi là quá muộn để khiến mọi người xung quanh nghi ngờ.

Onee-sama liếc nhìn tôi, đôi mắt nghiêm mà dịu: ``Reiko, lần này em không được tách ra. Nhớ không?''

Tôi đỏ mặt một chút, cúi đầu lí nhí: ``V-Vâng ạ\dots''

Con đường xuống núi dần hiện ra cảnh sắc chiều muộn: ánh vàng loang loáng lọt qua hàng cây, nhuộm cả những bụi bỉ ngạn đỏ thành sắc rực rỡ chói mắt. Tôi thoáng rùng mình, bước nhanh hơn để bám theo hai người.

Khi chạm chân xuống bậc đá cuối cùng, tôi không kìm được, quay đầu nhìn lại. Ngôi đền trên cao đã mờ đi sau lớp sương nhạt, chỉ còn lại cánh cổng torii đỏ nổi bật. Trong lòng tôi thoáng một cảm giác khó tả, vừa như bị thôi thúc, vừa như bị cảnh báo.

Tôi siết chặt quai túi, rồi bước lên xe buýt cùng hai anh chị. Cánh cửa khép lại, bánh xe lăn đi, để lại phía sau bóng đền chìm dần vào chiều sẫm.

Nhưng ngay khi xe vừa chuyển bánh, tôi chợt thoáng thấy một bóng hình lạ lẫm đang leo lên những bậc thang đá dẫn tới ngôi đền.

Một cô gái với đôi tai thỏ trắng muốt nổi bật trên mái tóc xanh dương nhạt, tay cầm chiếc máy ảnh kỹ thuật số hiện đại. Khác hẳn với Saya-san, Kaoru-san và Miku-san - những bóng ma với ánh mắt trống rỗng và cử chỉ máy móc - cô gái này toát lên một sức sống kỳ lạ. Mỗi bước chân của cô dường như mang theo một niềm vui thích thuần khiết, và ánh mắt cô sáng lên đầy nhiệt huyết khi ngắm nhìn khung cảnh xung quanh qua ống kính. Cô không hề có vẻ mờ ảo hay vô hồn như những bóng ma tôi từng gặp - ngược lại, cô hiện diện rõ ràng và chân thực đến kỳ lạ giữa không gian hoàng hôn mờ ảo này.

Rốt cuộc thì\dots\ cô ấy là ai?

\ct

Năm giờ chiều.

Xe buýt đưa chúng tôi về lại thị trấn khi bầu trời bắt đầu chuyển về màu đen, những màn sương đêm bắt đầu xuất hiện. Chúng tôi bước vào nhà, không gian lạ lẫm và hỗn độn đập vào mắt - một căn phòng khách chưa được sắp xếp gọn gàng, với đống hộp các-tông chất đống ở góc tường, những mảnh giấy gói và băng keo vương vãi khắp sàn.

Onee-sama ngả người xuống chiếc ghế sô-pha vừa được tháo khỏi lớp bọc ni-lông, thở dài: ``Trước khi trời tối hẳn, chúng ta nên nghỉ lấy sức.''

Tôi gật đầu, còn Onii-sama chỉ cười nhạt: ``Ừ, chứ không thì tối nay dễ thành gánh nặng lắm.''

Chúng tôi chia nhau mỗi người một góc trong căn phòng tạm bợ này. Xung quanh là những thùng đồ chưa mở - có cái ghi ``Đồ bếp'', có cái đề ``Sách vở'' - tất cả đều được chuyển đến từ nhà cũ, nơi bố mẹ vẫn đang bận rộn với thủ tục bàn giao và dọn dẹp. Chị tôi nằm dài trên ghế sô-pha, anh tôi ngồi dựa lưng vào bức tường trống trơn với cuốn sổ nhỏ, còn tôi cuộn mình trong chiếc chăn mỏng trên tấm đệm tạm được trải dưới sàn.

Tôi khẽ nhìn hai người họ qua ánh đèn vàng nhạt hắt xuống từ bóng điện trần trơ trọi - chị tôi với hơi thở đều đặn, anh tôi chăm chú ghi chép điều gì đó. Dù xung quanh vẫn còn bừa bộn với những hộp đồ chưa mở, những tấm bìa các-tông vứt vội một góc, nhưng bầu không khí bỗng dưng mang lại một chút ấm áp lạ kỳ. Như thể ba chúng tôi thật sự là một gia đình nhỏ, cùng nhau xây dựng tổ ấm mới này. Tôi khẽ mỉm cười, cảm nhận được sự gần gũi thân thương ấy, rồi nhắm mắt lại, để mặc cơ thể trôi vào giấc chợp mắt ngắn ngủi, trước khi đối mặt với cuộc phiêu lưu khó khăn phía trước.

\ct

Tám giờ ba mươi tối.

Chúng tôi rời nhà trong bóng tối, những bước đi có phần khẩn trương. Trời đêm lành lạnh, từng ánh đèn đường vàng hắt xuống, dẫn lối tới một quán ăn nhỏ nằm ở góc phố.

Khác với buổi trưa, quán này mang dáng vẻ cổ kính hơn, với ánh đèn mờ và hương súp thoang thoảng. Chủ quán là một bác trai lớn tuổi, tóc đã bạc quá nửa.

Chúng tôi ngồi xuống bàn gỗ cạnh cửa sổ, gọi vài món đơn giản. Vừa khi bác chủ rời đi, tôi bất chợt nghe thấy tiếng xì xào từ bàn bên.

``Nghe nói cô gái ở phố kế bên cũng biến mất rồi\dots''

``Thật sao? Cô ấy là ai vậy?''

``Là cô gái trẻ hay đi qua ga Aogami ấy. Nghe đâu tối qua đi làm về qua ga, rồi không thấy về nhà nữa.''

``Lại thêm một người nữa\dots\ Dạo này người mất tích nhiều quá.''

Tôi thoáng giật mình. Tôi liếc sang Onii-sama và Onee-sama, cả hai đều giật mình, ánh mắt giao nhau trong thoáng chốc.

Bác chủ quán quay lại, lắc đầu nặng nề khi nghe những lời đó: ``Lại thêm một vụ nữa à\dots\ mấy hôm nay chẳng yên gì cả.''

Tôi chợt nhớ về tờ báo mà tôi đã đọc từ hồi sáng, nơi mà đề cập đến số lượng người mất tích tăng đột biến và những tiếng còi tàu bí ẩn vang lên giữa đêm khuya.

Tôi ngập ngừng, đi về phía bác chủ quán, rồi lấy can đảm mở lời: ``Bác ơi\dots\ cháu có đọc được một bài báo, nói về số lượng người mất tích ở thị trấn. Thật sự là nhiều đến vậy ạ?''

Bác thở dài, chống tay lên quầy.
``Ừ, đúng thế. Cứ dăm bữa lại có người biến mất. Có người đi qua ga Aogami rồi chẳng bao giờ trở về. Người thì rời nhà buổi tối, không thấy sáng quay lại. Ai cũng sợ, nhưng dần dà rồi\dots\ thành quen. Thấy mất tích thì báo, cảnh sát tìm không ra, thế là thôi.''

Onii-sama ngả người ra sau, giọng thấp: ``Vậy còn\dots\ tiếng còi tàu thì sao? Có ai thực sự nghe thấy chưa ạ?''

Bác chau mày, nhớ lại, mặt gợn lên một sự mơ hồ. ``Tiếng đó à\dots\ có người nói nghe giống còi tàu lửa dùng than củi ngày xưa, cái loại cách đây mấy chục năm rồi không còn nữa. Bác thì chưa nghe, nhưng mấy cụ già sống gần ga, thỉnh thoảng bảo họ nghe thấy lúc nửa đêm. Mà không ai dám tới kiểm chứng.''

Lời kể ấy khiến không khí trên bàn trùng hẳn xuống. Tôi thoáng rùng mình, nhớ lại từng chữ trong bài báo buổi trưa. Những gì Kaoru-san nói hôm qua, hóa ra không phải lời bịa đặt, mà là lời đồn đã in sâu trong thị trấn này từ lâu.

Onee-sama khẽ nhíu mày, kết luận:
``Vậy là chắc chắn có gì đó không ổn ở ga Aogami. Chúng ta phải đi, nếu muốn tìm lời giải.''

Tôi gật đầu, lòng vừa run vừa háo hức. Chị liền siết chặt tay ly nước, khẽ nói: ``Được rồi. Ăn xong, chúng ta sẽ thẳng tiến tới ga. Không chần chừ gì thêm nữa.''

Trái tim tôi đập nhanh, nhưng tôi biết rằng, đêm nay chính là lúc chúng tôi phải đối mặt với sự thật.

\pov{Haragen Kyuen}

Chín giờ tối.

Chúng tôi rời quán ăn trong yên tĩnh, tiếng chuông gió trên cửa khẽ leng keng khi cánh cửa khép lại. Ngoài kia, thị trấn đã chìm hẳn vào đêm. Những ngọn đèn đường vàng vọt hắt xuống vỉa hè ẩm ướt, gió đông rít qua từng ngõ nhỏ, khiến con đường chúng tôi đi lúc ban ngày bỗng trở nên xa lạ.

Tôi bước đi ở giữa, hai bên là Saikyou và Reiko. Cả ba không ai lên tiếng ngay. Chỉ có tiếng bước chân hòa cùng nhịp thở, kéo dài trên mặt đường.

Cuối cùng, cô em lớn khẽ mở lời: ``Những gì ta gom lại từ sáng tới giờ\dots\ tiếng còi tàu, những vụ mất tích, việc Akane-san và Yuka-san không tới, rồi cả cục pin lượm được ở ngôi đền, và cả lời dặn của Shino-san\dots\ tất cả đều chỉ về ga Aogami.''

Tôi gật đầu.

``Ừ. Dù chưa thấy sợi dây liên kết nào rõ ràng, nhưng cái cách chúng xoay quanh cùng một chỗ\dots\ thì chắc chắn không phải là mấy thứ rời rạc ngẫu nhiên.''

Reiko ôm chiếc túi xách, giọng nhỏ nhẹ: ``Nhưng mà\dots\ em vẫn thấy sợ. Giống như là lần trước. Chúng ta cũng đã nói chuyện với nhau y như thế này, trên đường tới trường\dots\ lúc mà Shino-san và các bạn học của cậu ấy\dots''

Tôi thoáng khựng lại. Không cần để em ấy dứt câu, tôi cũng tờ mờ nhớ được. Vòng lặp trước - trận thử thách gan dạ do chính Honoka-san bày ra - cái kết ám ảnh ấy vẫn in hằn trong trí nhớ. Shino-san và những bạn học của cô ấy, kể cả người bày ra trò chơi đó\dots\ tất cả đều bỏ mạng. Sakura-san thì biến mất, còn tôi cùng Saikyou thì bị thương nặng. Tôi vẫn nhớ cái cảm giác máu loang trên tay mình, cái lạnh buốt len vào xương.

Tôi hít một hơi, cố giữ giọng vững vàng: ``Nhưng lần này thì khác. Chúng ta đã biết điều gì đang chờ phía trước. Và quan trọng hơn, ta vẫn còn ở cạnh nhau. Vẫn còn có sức để đi, để nói chuyện, chứ không giống như lần đó.''

Saikyou khẽ cười, đôi mắt dịu đi phần nào. Reiko cũng gật đầu, dù gương mặt vẫn phảng phất nét lo lắng. Khoảnh khắc ấy, tôi thấy cả ba như một gia đình thực thụ, cùng san sẻ nỗi sợ, cùng nương tựa để bước tiếp.

Chúng tôi tiếp tục đi. Đúng lúc đó, từ phía con hẻm bên phải, một bóng người vụt qua. Tôi dừng lại.

Một cô gái, mái tóc xanh dương nhạt xõa dài, đôi tai thỏ trắng nổi bật trong ánh đèn mờ, đang bước vội qua ngã rẽ. Khoảnh khắc ấy, tôi có cảm giác khác thường, không phải kiểu vô hồn đáng sợ, mà\dots\ một sự sống thật sự. Một người bình thường.

Tôi định gọi với theo, nhưng bóng dáng ấy đã biến mất sau khúc cua, để lại khoảng trống lạnh lẽo và cơn gió thổi tạt qua má.

Reiko nhận ra tôi dừng lại, khẽ hỏi:
``Onii-sama? Có chuyện gì vậy?''

Tôi lắc đầu, nuốt chậm một hơi.
``Không\dots\ chỉ là nhìn nhầm thôi.''

Dù tôi biết, mình đã không hề nhầm.

Chúng tôi tiếp tục đi, từng bước chậm dần trên con đường nhỏ. Bầu không khí nặng nề, mùi ẩm mốc từ những ngõ tối, tiếng gió lùa khe khẽ. Tôi nhìn về phía trước, nơi bóng đêm như dày đặc hơn, đen kịt hơn, phả dới lớp sương đêm.

Có gì đó đang chờ đợi ở ga Aogami. Một điều mà cả ba chúng tôi chưa từng tưởng tượng nổi.

Một cái gì đó bí ẩn? Kinh hoàng? Đáng sợ? Ám ảnh?

Chúng tôi không thể nào biết được. Trăm nghe không bằng thấy, trăm thấy không bằng một chạm. Đến giờ phút này\dots

\dots tự tay chúng tôi sẽ điều tra về sự bí ẩn trong ga tàu kia. Tìm ra Yuka-san và Akane-san thực sự đang ở đâu.

\end{document}